\section{Shared calculus for local moments}\label{app:local-calculus}\label{app:catalogue-map}\label{app:local-catalogue}
This appendix collects the scalar tools shared by the local estimates and
records where each one is used. In every statement, the direction and the
weight are part of the hypotheses. References to Appendices~\ref{sec:three}--\ref{app:five-remaining-outlets} point to
the technical appendix that accompanies this article.
\begin{table}[htbp]
\centering
\caption{Local estimates and their uses in the assembly.}\label{tab:proof-index}

\begin{tabular}{@{}p{30mm}p{74mm}p{43mm}@{}}\toprule
Location & Estimates & Used in\\\midrule
Appendix~\ref{app:path-moments} & Path fourth moments; singleton fifth and cubed weights & terminal and center bounds\\
Appendix~\ref{app:zero-moments} & Unpropagated even and odd messages & zero-step supply; canonical roots\\
Appendix~\ref{app:even-propagation} & Even residual lists after positive propagation & terminal supply\\
Appendix~\ref{app:odd-propagation} & Odd residual lists after positive propagation & terminal supply\\
Appendix~\ref{app:center-moments} & Cubed-weight estimates for nonsingleton centers & four-branch star assembly\\
Appendix~\ref{app:five-directions} & Index to direction tables and short-length exceptions & all five-branch constructions\\
Appendices~\ref{app:five-local-proofs}--\ref{app:five-remaining-outlets} & Reference fourths, ordinary eighths and common-reference outlets & five-branch boundary estimates\\\bottomrule
\end{tabular}

\end{table}
The common spectral and rooted-spider foundations are in
Appendix~\ref{q3:app:foundations}.\label{loc:S1}
The scalar tools below carry no arm parameters; each application supplies
its own messages, signed means and parameter interval. A finite
coefficient table records a polynomial identity or a bound on a
continuous domain.

\paragraph{Common scalar calculations.}
For a nonnegative weight $w$ with $\E w=1$, write
$\E_w H=\E[wH]$; the measure and the adjoint of $K$ are unchanged.
Each application specifies its weight and message.

If $|R|\le M$, then, for $j\ge2$,
$|\E_w R^j|\le M^{j-2}\E_wR^2$. In particular, set
\begin{equation}\label{eq:compact-binomial}
 B_r(m,M)=\sum_{j=2}^{r}\binom rj
                 m^{r-j}M^{j-2}2^{-j}\qquad(r=6,8,12).
\end{equation}
For $m\ge0$ and $0\le P\le(m+R/2)^2$, binomial expansion gives
\begin{equation}\label{eq:compact-power-bound}
 \E_w P^{r/2}\le m^r+\frac r2m^{r-1}\E_wR
                  +B_r(m,M)\E_wR^2 .
\end{equation}
Every coefficient of $B_r$ is nonnegative. Thus upper bounds for $m,M$
can be substituted term by term, retaining the signed linear mean separately.
The fourth-power expansions below specify their center and weight.

If instead $P\le\exp(R-\Delta)$, let
$\Lambda(t)=\log\E_w e^{tR}$. Its second derivative is the variance of
$R$ under the tilted probability measure, and is at most $M^2$. Hence
\begin{equation}\label{eq:compact-exponential}
 \E_w P^s\le
 \exp\{s(\E_wR-\Delta)+s^2M^2/2\}\qquad(s\ge0).
\end{equation}
In particular $\Delta-\E_wR\ge3M^2$ suffices for the sixth moment.
For the centered square envelopes, Cauchy--Schwarz also gives
\begin{equation}\label{eq:compact-weighted-variance}
 V=\E R^2,\quad \|R\|_\infty\le M
 \quad\Longrightarrow\quad
 \E[pR^2],\ \E[qR^2]\le V+M\sqrt{dV}.
\end{equation}

\paragraph{Reading the interval tables.}
The tables below specify upper bounds on closed intervals. When an upper
bound has the form $H(d)+ad$ and $H''\ge0$, the two listed endpoint bounds
prove the assertion throughout that interval. A decreasing bound needs only
the left endpoint. Entries are rational bounds from the displayed finite
polynomials or their geometric tails; intermediate integer cross-products
are suppressed.
The case $d=0$ is settled by leaf integration: every actual rooted message
equals one, so divisions by $d$ below are confined to $d>0$. Equal-pair
restoration and complete reversal follow Section~\ref{sec:prelim}.

\subsection{Scalar and propagation tools}\label{calc:logtools}
The scalar inequalities used for odd-root and propagation estimates are
\begin{equation}\label{eq:compact-log-tools}
 \begin{gathered}
 \log(1-x)\le-x-3x^2/10\quad(-1\le x<1),\\
 \log(1+x)\le x-3x^2/10\quad(-1<x\le1),\\
 e^{y-y^2/7}\le1+y+3y^2/8,\qquad e^{y-y^2/8}\le1+y+5y^2/12.
 \end{gathered}
\end{equation}
For the logarithm, differentiate
$\log(1-x)+x+3x^2/10$; its derivative is $-x(2+3x)/[5(1-x)]$,
and its only possible maxima are $-1,0$, with $\log2<7/10$.
For the last two inequalities the positive quadratic polynomials $Q_7,Q_8$
have logarithmic differences with derivatives
$y(2-5y+6y^2)/(56Q_7)$ and $y(4-8y+5y^2)/(48Q_8)$,
respectively. Each quadratic factor is positive, proving the global bounds.

We also need the sharper global bound, for every real $x$,
\begin{equation}\label{eq:shared-exp-eighth}
 e^{x-x^2/8}\le Q(x):=1+x+13x^2/32,
\end{equation}
because $[\log Q-x+x^2/8]'=x(13x^2-20x+8)/(128Q)$ has the sign of $x$.

\paragraph{Shared propagation identities.}\label{calc:duhamel}
Write $\operatorname{Alt}_h=\sum_{i=1}^h(-1)^{i-1}\delta_{i}$ for each finite $h\ge0$.
For the propagation catalogue write $\tau=d-\delta_{2}$ and $S'_n=\sum_{j=0}^nK'^j\one$.
Here and below $K'=-K$, $g'=-g$, $\bar f_{n}=T^n\one$ and
$\|F\|_{4,p}=(\E[p|F|^4])^{1/4}$.
The finite Duhamel identity, obtained by iterating $T=J+K'$, is
\begin{equation}\label{eq:compact-duhamel}
 T^Lu=K'^Lu+\sum_{j=0}^{L-1}\E[K'^ju]\bar f_{L-1-j}.
\end{equation}
The actual path means also satisfy the stronger uniform loss
\begin{equation}\label{eq:compact-strong-path-means}
 \mu_n=\E f_n=\E \bar f_{n}\le1-83d/32\qquad(n\ge4).
\end{equation}
For $n=4,5,6$ use their exact formulas; the losses divided by $d$ are at
least $83/32,137/48,89/32$. For
$\mu_7=1-6d+4\delta_{2}+6d^2-2\delta_{3}-2d\delta_{2}$ the loss is at least $127/48$.
For $n\ge8$ retain $f_0,\ldots,f_5$ in convolution. The three preceding
beta indices are at least seven, each contributing an additional
$c_0\delta_{3}$, $c_0=19/32$. Hence
\[
 1-\mu_n\ge\tau(7-10d+3\delta_{2}+d^2)+3c_0\delta_{3}
 \ge c_0d(7-10d+7d^2)\ge95d/36>83d/32.
\]
This uses integral moment inequalities $\delta_{2}\ge d^2$, $\delta_{3}\ge d^3$,
without imposing a lower bound on individual spectral frequencies.

\paragraph{A finite perturbation of $G_n$.}
If a spectral polynomial has the form $G_n(z)+E_n(z)$ with
\[
 |E_n(i\sqrt x)|\le(4/5)x^{(n-1)/2},\qquad 0\le x\le\theta,
\]
then its modulus is at most one for every $n\ge4$.
For odd $n\ge5$, use
$1-|G_n|\ge(x-x^n)/[2(1+x)]$; the required worst-index comparison is
$17/18+625/20736=20209/20736<1$.
For $n=4$, $|G_4|\le1-x/2-x^2/2$ and
$(4/5)x^{3/2}\le x(1+x)/2$ by a positive quadratic in $\sqrt x$.
For $n\equiv0\pmod4$, $n\ge8$, compare against
$1-1/\sqrt{1+x}\ge x/[2(1+x)]$; the worst bound is
$3/(5\sqrt2)<1$.
For $n\equiv2\pmod4$, $n\ge6$, the sufficient inequality is
\[
 (8/5)(1+x)x^{(n-3)/2}+2x^{n/2-1}+x^{n-1}\le1.
\]
At the worst index $n=6$, use $x\le5/12$, $\sqrt x\le2/3$ and
$x^5<1/64$ to obtain $17/27+25/72+1/64=1715/1728<1$.
All higher powers improve these bounds; at $x=0$ there is equality.

\paragraph{Fourth-moment endpoint comparison.}
For $Z\ge0$, put $\mu=\E Z$, $w=Z-\mu$ and $h=\E[-gw]$.
The original $p$-weighted expansion is exactly
\[
 \E[pZ^4]=\mu^4-4\mu^3h+
            \E[pw^2(Z^2+2\mu Z+3\mu^2)].
\]
If $0\le\mu\le\bar\mu$, $h\ge h_0\ge0$,
$\operatorname{Var}Z\le V$ and the pointwise coefficient is at most $M$,
then
\begin{equation}\label{eq:compact-fourth-endpoints}
 \E[pZ^4]\le MV+\max\{0,\bar\mu^3(\bar\mu-4h_0)\}.
\end{equation}
This treats the two opposite-sign mean terms together: the maximum of
$u^3(u-4h_0)$ on $[0,\bar\mu]$ is at an endpoint.


For later error estimates,
\begin{equation}\label{eq:compact-qnorm}
 \E[pq^4]\le1+12d/5,\qquad \|q\|_{4,p}\le e^{3d/5}.
\end{equation}
Indeed $pq^4=1-3g+2g^2+g^3(1-g)(2-g)$; its last term is nonpositive
when $g\le0$, and for $g\ge0$ use
$g(1-g)(2-g)\le2/(3\sqrt3)<2/5$.
The elementary exponential remainder used with these bounds is
\begin{equation}\label{eq:compact-exp-remainder}
 e^x\le1+x+\frac{x^2}{2(1-x/3)}\qquad(0\le x<3).
\end{equation}
Its proof bounds $j!\ge2\,3^{j-2}$ for every $j\ge2$ and sums the
positive geometric tail.

\subsection{A fourth-power coefficient}\label{calc:pathscalar}
Set $c=19/32$, $x=\sqrt d$, $x_*=2/\sqrt{19}$,
$m=1-3cx^2$, $b_0=1-cx^2$ and $a_0=1+m+x$.
For $B_x(q)=a_0-b_0q$ put
\begin{equation}\label{eq:compact-pathscalar}
 H_x(q)=q[B_x(q)^2+2mB_x(q)+3m^2]\le7-10x^2
 \quad(0\le x\le x_*,\ 0\le q\le2).
\end{equation}
Here $B_x\ge0$, and
$H_x''=6b_0^2q-4b_0(a_0+m)\le4b_0x(3cx-1)\le0$.
At $q=3/2$ the value and derivative are
\begin{align*}
 H_0(x)&=51m^2/8+9mx/2+3x^2/2,\\
 D(x)&=-1/4-221x^2/64-57x^3/16+35739x^4/4096.
\end{align*}
The derivative $(-D)'=x[221/32+171x/16-35739x^2/1024]$ is positive:
the bracket is concave, and its endpoint lower bound at $x_*$ is
$93233/24320>0$. Thus $D\le-1/4$ and the maximum in $q$ lies in
$[0,3/2]$. On that interval $H_x''\le-1813/704<-5/2$, using
$b_0\ge7/8$ and the endpoint bound
$-3-4x+15cx^2\le-259/88$.
The tangent bound is therefore $\max H_x\le H_0+D^2/5$.
The two derivative certificates are $-D(1/3)<2/3$ and
$-D(x_*)<89205/92416<1$.
Write
\[
 Q=H_0+10x^2-7
 =-5/8+9x/2-1435x^2/128-513x^3/64+165699x^4/8192.
\]
For $x\le1/3$ its last two terms have nonpositive sum, so
$Q+D^2/5<-5/8+1/2+4/45=-13/360$.
For $1/3\le x\le x_*$ its quadratic part is at most $-427/1152$;
the last two terms sum to less than $2/15$, using $x_*<6/13$.
Hence $Q+D^2/5<-427/1152+2/15+1/5=-43/1152$.
This proves the full continuous domain of \eqref{eq:compact-pathscalar}.

\subsection{A common fourth-moment comparison}\label{calc:fourthcomparison}
Set $c=19/32$, $\tau=d-\delta_{2}$ and $g'=-g$. Suppose $Z=S-a-bq\ge0$, with
\begin{gather*}
 \E S=1-\alpha,\quad\E[g'S]=\eta,\quad
 a,b,\alpha,\eta\in[\tau,d],\quad\|S-1\|^2\le d,\\
 |S-1-g'|\le\sqrt d.
\end{gather*}
For $\mu=\E Z$, $w=Z-\mu$ the exact identities imply
\begin{gather*}
 0\le\mu\le m=1-3cd,\quad \E[g'w]\ge h=cd-d^2,\\
 \E w^2=\|S-1\|^2-\alpha^2+b^2d-2b\eta
                  \le d-3c^2d^2+d^3,\\
 Z\le(1-cd)q-cd+\sqrt d.
\end{gather*}
More generally, the four displayed output bounds on the mean, covariance,
variance and pointwise envelope alone suffice for the conclusion below;
the coefficient conditions $a,b,\alpha,\eta\in[\tau,d]$ are one
sufficient route to those outputs. In particular, later applications
may have $b>d$ if they verify the outputs directly.
The coefficient $p(Z^2+2\mu Z+3\mu^2)$ is bounded by $7-10d$ on the
whole $0\le d\le1/3$. For $d\le4/19$ this is
\eqref{eq:compact-pathscalar}. For the extension put $x=\sqrt d$ and use
the same polynomials $H_0,D,Q$ from its proof. On
$[2/\sqrt{19},1/\sqrt3]$, $-D$ increases to less than $7/6$,
$H'_x(2)<0$, and $H''_x\le-3/2$ on $[0,3/2]$.
The latter follows from $b_0\ge77/96$ and
$-3-4x+15cx^2\le-519/224$.
Also $Q'<0$: it is convex and has negative endpoint values
\[
 Q'(2/\sqrt{19})=-171/304-(52421/9728)(2/\sqrt{19}),\quad
 Q'(1/\sqrt3)=-225/64+(9313/2048)/\sqrt3.
\]
Since $Q(2/\sqrt{19})<-3/4$, the tangent remainder is at most
$-3/4+49/108=-8/27<0$.
In the fourth expansion the combined function
$\Phi(u)=u^3(u-4h)$ is maximized on $[0,m]$ at $m$:
$m-4h=1-7cd+4d^2\ge17/288>0$.
It follows that
\begin{equation}\label{eq:compact-fourth-comparison}
 \E[pZ^4]\le R_0:=m^3(1-7cd+4d^2)+(7-10d)(d-3c^2d^2+d^3)
 \le1-d/9.
\end{equation}
The divided coefficients of $R_0-1$ are bounded by
$-5/2,37/2,-49,52,-22$. The resulting polynomial is at most
$-5/2+(37/2)d-36d^2\le-71/576$ for $d\le1/4$.
For $d=1/4+h$, $h\le1/12$, its exact expansion is
$-27/128+19h/8-73h^2/4+30h^3-22h^4$;
using $30h^3\le5h^2/2$ bounds it by $-979/8064<-1/9$.

\proofpart{An enlarged short tail}
The same comparison holds if the row tail is $(13/12)\sqrt d$ and
$d\le1/10$. At $p=3/2$ the modified derivative and remainder are
\begin{align*}
 D_\lambda&=-1/4-1889x^2/576-247x^3/64+35739x^4/4096,\\
 Q_\lambda&=-5/8+39x/8-4205x^2/384-2223x^3/256+165699x^4/8192.
\end{align*}
Here $|D_\lambda|\le2/3$, the maximum is in $p\le3/2$, and
$H''\le-14/5$. The cubic and quartic sum in $Q_\lambda$ is nonpositive;
the remaining quadratic and tangent gain leave
$-5/8+4563/8410+5/63=-6499/2119320<0$.

\subsection{Five short-connector coefficient profiles}\label{calc:coefficientprofiles}
For each row below put $x=\sqrt d$, $u$ equal to its mean bound,
$B$ equal to its row envelope, and $H(\xi)=\xi(B^2+2uB+3u^2)$, with $q=2-\xi$.
All rows have $H\le7$. This follows from concavity, $H'(3/2)<0$,
and the tangent estimate $H\le H(3/2)+D^2/(2\kappa)$ when
$H''\le-\kappa$ below $3/2$.
The displayed formulas define every coefficient in the table; bounds
are on $D=-|D|$ and $Q=H(3/2)-7$.
\begin{table}[htbp]
\centering
\caption{Scalar coefficient profiles for short connectors.}\label{tab:profiles}
\small
\begin{tabular}{@{}cllll@{}}
\toprule
 & $(u,B)$ & range of $x$ & $(|D|,\kappa)$ & upper bound for $Q$\\\midrule
A & $(1-d,(1-d/3)q-d/3+x)$ & all & $(9/10,8/3)$ & $-7/40$\\
B & $(1-3d/2,(1-d/2)q-d/2+13x/12)$ & all & $(11/10,5/2)$ & $-53/200$\\
C & $(1-9d/8,(1-3d/10)q+x)$ & $x\le1/3$ & $(21/40,27/10)$ & $-1/8$\\
 & & $x\ge1/3$ & $(1,27/10)$ & $-2/5$\\
D & $(1-5d/8,q+4x/5-3d/8)$ & all & $(1,3)$ & $-13/72$\\
E & $(1-d,q+x-2d/3)$ & $x\le1/3$ & $(2/3,3)$ & $-13/64$\\
 & & $x\ge1/3$ & $(7/5,3)$ & $-1/2$\\\bottomrule
\end{tabular}

\end{table}
Differentiation gives
$H''=6b^2\xi-4b(A+u)$ when $B=A-b\xi$.
In each row $-D$ increases; its derivative divided by $x$ is an
endpoint-positive concave quadratic. On the upper intervals of C and E,
$Q'$ is convex with negative endpoints, giving the tabulated bounds;
for A, D and the lower intervals of C and E, discard the nonpositive
cubic/quartic sum and maximize the remaining quadratic.
For B instead use $x\le3/5$ to bound that sum by $(207/160)x^3$;
then $Q_B\le-5/8+39x/8-(9953/600)x^2<-53/200$.
\subsection{A positive cycle form}\label{loc:S6}
Let $r=Kg$, $m_j=\E g^j$ and
$C(x,y,z)=(1+K(x,y))(1+K(y,z))(1+K(z,x))\ge0$.
For every real $f$ the following exact form is nonnegative:
\begin{equation}\label{eq:compact-cycle-form}
 \E[(1+2r)f^2]+(\E f)^2-\|Kf+gf\|_2^2
 =\tfrac12\E[C(x,y,z)(f(x)+f(y))^2]\ge0.
\end{equation}
To verify it, integrate $C$ first in one or two coordinates:
\begin{align*}
 \E_{y,z}C(x,y,z)&=1-g(x)^2+2r(x),\\
 \E_zC(x,y,z)&=(1+K(x,y))[1+g(y)-g(x)+K^2(x,y)].
\end{align*}
The symmetric part of the second kernel is
$1+K^2(x,y)+K(x,y)(g(y)-g(x))$. Thus
\[
 \E[Cf(x)f(y)]=(\E f)^2-\|Kf\|_2^2-2\E[gfKf].
\]
Combining this with the first marginal proves
\eqref{eq:compact-cycle-form}; in particular $\E C=1-3d$.
Taking $f=g+t$ gives, with $\Delta=d-\delta_{2}-m_4$,
\[
 \Delta-4tm_3+2(1-3d)t^2\ge0\quad(t\in\mathbb R),
 \qquad \Delta\ge0,\qquad 2m_3^2\le(1-3d)\Delta.
\]
The last inequality also follows without division by Cauchy--Schwarz on
the nonnegative measure $C\,d\mu^3$: its quadratic form has values
$Q(g)=\Delta$, $Q(1)=2(1-3d)$, and $Q(g,1)=-2m_3$.
Repeated atoms are retained. Their diagonal contribution gives
$1-3d\ge\sum_x\mu(x)^3>0$ and
$\Delta\ge2\sum_x\mu(x)^3g(x)^2$, so $\Delta>0$ when $d>0$.
The quantities $\delta_{2},m_3,m_4$ satisfy the joint budgets above.

\subsection{Integer reference powers}\label{loc:S10}
Put $c=pq+2Kg\ge0$ and
$\sigma_m(z)=\sum_{j=0}^{m-1}2^{m-1-j}z^j$, with $\sigma_0=0$.
For every integer $m\ge0$,
\begin{equation}\label{eq:compact-reference-integer}
 A(p^m)\le\frac{p^{m+1}}{m+1}+\frac c2\sigma_m(p),\qquad
 T(q^m)\le\frac{q^{m+1}}{m+1}+\frac c2\sigma_m(q).
\end{equation}
These follow from the convex transport proved below: take $\phi(t)=t^m$
and use $(2^m-p^m)/q=\sigma_m(p)$. The cases $m=0,1$ are equalities;
in particular $f_2=(p^2+c)/2$ and $A(p^2)\le p^3/3+(p+2)c/2$.
Multiplication by a specified nonnegative function and integration preserve
these pointwise inequalities.

\subsection{Convex reference transport}\label{loc:S11}
For every continuous convex $\phi:[0,2]\to\mathbb R$,
\begin{equation}\label{eq:compact-convex-transport}
 A[\phi(p)](x)\le\int_0^{p(x)}\phi(t)\,dt+
 \frac{c(x)}{2q(x)}[\phi(2)-\phi(p(x))].
\end{equation}
Complete reversal gives the corresponding $T[\phi(q)]$ bound.
Affine tests give equality.

First, if $s$ is the score of any finite skew kernel bounded in absolute value by one and $P=1+s$, then
\[
 \E(P-a)_+\le(2-a)^2/4\quad(0\le a\le2),\qquad
 \E\psi(P)\le\tfrac12\int_0^2\psi(t)\,dt
\]
for every continuous convex $\psi$. Indeed, for $S=\{P>a\}$ and
$z=\mu(S)$, skewness gives $\int_Ss\,d\mu\le z(1-z)$, and hence
$\E(P-a)_+\le z(2-a-z)\le(2-a)^2/4$.
Express a convex piecewise linear function as an affine function plus
a nonnegative combination of $(t-a)_+$, then take uniform polygonal
approximations. The constants and linear terms agree because $\E P=1$.
In particular, for a scalar $0<P_0\le2$,
\begin{equation}\label{eq:compact-scaled-score}
 \E\phi(P_0P/2)\le P_0^{-1}\int_0^{P_0}\phi(t)\,dt.
\end{equation}

Write $B=T^2$. Its kernel is
$B(x,y)=1-g(x)+g(y)+K^2(x,y)$; nonnegative two-step convolution gives
$0\le B(x,y)\le2q(x)$. The cycle marginal just proved is
\[
 \E_y[A(x,y)B(x,y)]=c(x)\ge0.
\]
At a fixed $x$, set $p_x=p(x)>0$. For this scalar comparison alone use
$d\nu_x(y)=A(x,y)d\mu(y)/p_x$. Its score is
$\gamma_x(y)=[g(y)-K^2(x,y)]/p_x$. Put
$a_y=p_x(1+\gamma_x(y))/2$, $b_y=B(x,y)/2$. Then
\[
 0\le a_y\le p_x,\quad0\le b_y\le q_x,\quad
 a_y+b_y=p(y),\quad \E_y[A(x,y)b_y]=c(x)/2.
\]
For $0\le a\le p$, $0\le b\le q$, convexity gives
\[
 \phi(a+b)-\phi(a)\le\frac bq[\phi(a+q)-\phi(a)]
 \le\frac bq[\phi(2)-\phi(p)].
\]
The second inequality is the monotonicity of a convex increment of fixed
length, also valid for negative increments. Integrate against
$A(x,y)d\mu(y)$ and apply \eqref{eq:compact-scaled-score} to $\nu_x$.
This proves \eqref{eq:compact-convex-transport}; all target operators and
adjoints still use the original $\mu$.

For $\alpha\ge1$, the tests $t^\alpha$ and $(2-t)^\alpha$ give
\begin{align*}
 A(p^\alpha)&\le\frac{p^{\alpha+1}}{\alpha+1}
                  +\frac{c(2^\alpha-p^\alpha)}{2q},\\
 A(q^\alpha)&\le\frac{2^{\alpha+1}-q^{\alpha+1}}{\alpha+1}
                  -\frac c2q^{\alpha-1}.
\end{align*}
Their full reversals give the other two real-power inequalities.
The test $(t-h)_+$, $0\le h\le2$, gives
\[
 A[(p-h)_+](x)\le\frac{(p(x)-h)_+^2}{2}
 +\frac{c(x)}{2q(x)}[(2-h)-(p(x)-h)_+].
\]
The negative cycle term in the opposite-direction power bound is retained.



\subsection{Constants and arithmetic certificates}\label{sec:certificate-scope}
\begin{table}[htbp]\centering\small
\caption{Fixed constants and their roles. The rational relaxations are sufficient bounds.}\label{tab:constants}
\begin{tabular}{@{}p{26mm}p{55mm}p{65mm}@{}}\toprule
Constant & Origin & Use\\\midrule
$2/\pi$ & Skew-kernel norm bound, Appendix~\ref{q3:app:foundations} & Spectral estimates.\\
$\theta=13/32$ & Rational upper bound for $(2/\pi)^2$ & Geometric control of $\delta_{j+1}/\delta_j$.\\
$\alpha=19/32$ & $1-\theta$ & Centered resolvent and fourth-moment losses.\\
$\rho=5/12$ & Convenient relaxation of $\theta$ & Geometric tails in Appendix~\ref{q3:app:operators}.\\
$C_*=101/50$ & Relaxed spectral mask bound in the four-branch estimates & Pair and multiple-arm envelopes.\\
$C_*=32768/16245$ & Sharper finite-mask bound used in the five-branch estimates & Joint terminal and root-mask estimates.\\
$\eta_*=1/500$ & Common sufficient outlet loss, Appendix~\ref{app:five-directions} & Two compatible outlets give loss $d/2000$.\\
$7/15$ & Residue-class geometric deficit, \eqref{eq:compact-G-loss} & Finite spectral polynomials.\\
$625$ & Negative-part estimate, Section~\ref{loc:S35} & Signed odd-table auxiliaries.\\\bottomrule
\end{tabular}
\end{table}
Each certificate follows the same pattern: a nonnegative message
envelope, a spectral polynomial, signed mean and variance bounds, and a
scalar moment comparison. Rational endpoints keep the last stage finite
while covering the whole continuous parameter domain.
Appendix~\ref{app:worked} works out one example in full.


The full local arguments are in the technical appendix.

\subsection{Polynomial coefficient certificates}
\label{sec:bernstein-certificates}

The polynomial comparisons below use rational coefficients on closed
intervals or rectangles. The following elementary criterion makes
their finite coefficient tables sufficient for the whole domain.

\begin{lemma}[Bernstein coefficient bounds]\label{lem:bernstein-certificate}
Let \(P\) be a polynomial on
\([a,b]\times[c,d]\), and write
\[
 Q(s,t)=P(a+(b-a)s,c+(d-c)t)
       =\sum_{p=0}^{n}\sum_{q=0}^{m}A_{pq}s^pt^q.
\]
Define
\begin{equation}\label{eq:bernstein-coefficients}
 B_{ij}=\sum_{p=0}^{i}\sum_{q=0}^{j}
 A_{pq}\frac{\binom{i}{p}}{\binom{n}{p}}
        \frac{\binom{j}{q}}{\binom{m}{q}},
 \qquad 0\le i\le n,\quad0\le j\le m.
\end{equation}
Then
\begin{equation}\label{eq:bernstein-convex-bound}
 \min_{i,j} B_{ij}\ \le P(x,y)\le\ \max_{i,j}B_{ij}
 \qquad ((x,y)\in[a,b]\times[c,d]).
\end{equation}
For a polynomial on one interval, take \(m=0\) and omit the
second variable. Degrees \(n,m\) may exceed the actual degrees.
The coefficients and both bounds are rational when the data are rational.
\end{lemma}
\begin{proof}
The binomial theorem gives
\[
 s^p=\sum_{i=p}^{n}
       \frac{\binom{i}{p}}{\binom{n}{p}}
       \binom{n}{i}s^i(1-s)^{n-i}.
\]
Apply this identity to both variables. It yields
\[
 Q(s,t)=\sum_{i=0}^{n}\sum_{j=0}^{m}B_{ij}
 \binom{n}{i}s^i(1-s)^{n-i}
 \binom{m}{j}t^j(1-t)^{m-j}.
\]
For \(s,t\in[0,1]\), the weights are nonnegative and sum to one.
This proves \eqref{eq:bernstein-convex-bound}, including the boundary.
\end{proof}


\subsection{Two reference curvature bounds}
\label{sec:reference-curvature-nine}

\begin{lemma}\label{lem:reference-coefficient-nine}
Let \(0\le d\le1/3\), \(t=\sqrt d\), and let \(q,C,Z\) be
nonnegative functions on a probability space, with \(q\le2\)
and \(\mu=\E Z\). Define the following scalar functions:
\begin{equation}\label{eq:reference-coefficient-nine:data}
\begin{aligned}
 \mathcal C_w&=2-q+6t/5,&
 \mathcal M_w&=1-3t^2/2,\\
 \mathcal S_w&=q+1+13t/12-(2+5q/8)t^2,\\
 \mathcal C_s&=2-q+6t/5-19t^2/32,&
 \mathcal M_s&=1-11t^2/8,\\
 \mathcal S_s&=q+1+t-(13/8+3q/8)t^2.
\end{aligned}
\end{equation}
If, for one \(i\in\{w,s\}\),
\[
 C\le\mathcal C_i,\qquad Z+\mu\le\mathcal S_i,
 \qquad \mu\le\mathcal M_i,
\]
then
\begin{equation}\label{eq:reference-coefficient-nine:bound}
 C\{(Z+\mu)^2+2\mu^2\}\le9.
\end{equation}
\end{lemma}
\begin{proof}
Enlarge the scalar domain to \(0\le t\le3/5\), \(0\le q\le2\).
Both \(\mathcal M_i\) are at least \(23/50\), and both
\(\mathcal S_i\) are at least \(93/100\): their coefficient of
\(q\) is positive, and their remaining concave quadratics have
these endpoint lower bounds. The two \(\mathcal C_i\) are
nonnegative as well, since \(6/5-19t/32>0\).
Thus the asserted coefficient is at most
\(\mathcal C_i(\mathcal S_i^2+2\mathcal M_i^2)\).

Put \(D_i=9-\mathcal C_i(\mathcal S_i^2+2\mathcal M_i^2)\).
Divide the enlarged domain at \(t=3/10\) and \(q=1\).
On each of the four rectangles, map both variables affinely to
\([0,1]\) and apply \eqref{eq:bernstein-coefficients}.
The complete coefficient lists have the following bounds:
\begin{table}[htbp]\centering\small
\caption{Complete coefficient bounds for the reference curvature comparisons.}
\label{tab:reference-coefficient-nine:bernstein}
\begin{tabular}{@{}ccccc@{}}
\toprule
Polynomial&Bidegree&Rectangles&Total coefficients&Lower bound\\
\midrule
\(D_w\)&\((5,3)\)&4&96&\(1/2\)\\
\(D_s\)&\((6,3)\)&4&112&\(1/2\)\\
\bottomrule
\end{tabular}
\end{table}

Lemma~\ref{lem:bernstein-certificate} gives \(D_i\ge1/2\)
throughout the enlarged domain, including its boundary.
This proves \eqref{eq:reference-coefficient-nine:bound} on
the original domain \(t\le1/\sqrt3<3/5\).
\end{proof}
